\section*{Capítulo \ref{ch:b1:fluid-envelopes-structure} --- Estructura de la atmósfera y del océano}
% values: calc/ch26.py

\begin{solution}{exo:b1:fluid-envelopes-structure:1}
(a)~$H = R^*T/(M_0 g_0) = 8314.32 \times 250/(28.9644 \times 9.80665) =
\qty{7.32}{km}$. (b)~$e^{10/7.32} = 3.92$: la presión baja hasta cerca de una cuarta parte.
(c)~$H$ es proporcional a la temperatura: a \qty{216.65}{K} vale
\qty{6.34}{km}, frente a \qty{8.43}{km} a \qty{288.15}{K}. El aire frío es más denso a una presión dada, así que la presión baja más deprisa con la altura.
\end{solution}

\begin{solution}{exo:b1:fluid-envelopes-structure:2}
Con $\Gamma_d = \qty{9.76}{K/km}$: \qtyrange{0}{0.3}{km}, $11.0 > 9.76$,
inestable; \qtyrange{0.3}{1.2}{km}, $9.7$, neutra dentro de la precisión de los datos; \qtyrange{1.2}{1.5}{km}, una inversión; \qtyrange{1.5}{8}{km}, estable;
por encima de \qty{8}{km}, isoterma, estable. La inversión es la más estable, pues $\Gamma_d - \Gamma = 13.8$ es allí el valor mayor.
\end{solution}

\begin{solution}{exo:b1:fluid-envelopes-structure:3}
(a)~$\sigma = 1026.95 - 1000 = \qty{26.95}{kg.m^{-3}}$. (b)~B, la más ligera, yace
encima de A. (c)~Todas las densidades del agua de mar empiezan por las mismas cifras, 10; $\sigma$ conserva solo las cifras que cambian de un agua a otra.
\end{solution}

\begin{solution}{exo:b1:fluid-envelopes-structure:4}
(a)~$S = 0.25 \times 35.40 + 0.75 \times 34.30 = 34.575$; $\theta = 0.25 \times
15.0 + 0.75 \times 4.0 = \qty{6.75}{\celsius}$. (b)~Sobre el segmento PQ, a un cuarto del camino de Q a P (\cref{prop:b1:fluid-envelopes-structure:mixing}).
(c)~Las mezclas de P, Q y R llenan el triángulo PQR. Escribiendo
$(34.60, 8.0) = a\,\mathrm P + b\,\mathrm Q + c\,\mathrm R$ con $a + b + c = 1$
se obtiene $a = 0.345$, $b = 0.788$, $c = -0.133$: una proporción negativa de R, así que el punto queda fuera del triángulo y ninguna mezcla de las tres puede producirlo.
\end{solution}

\begin{solution}{exo:b1:fluid-envelopes-structure:5}
(a)~$\theta = 253.15 \times 2^{2/7} = \qty{308.6}{K}$. (b)~A \qty{1000}{hPa},
$\theta = T = \qty{298.15}{K}$. La \omterm{def:b1:fluid-envelopes-structure:potential-temperature}{temperatura potencial} aumenta hacia arriba, de 298.2 a \qty{308.6}{K}: la capa es estable para aire seco. (c)~Llegaría al suelo con su \omterm{def:b1:fluid-envelopes-structure:potential-temperature}{temperatura potencial}, \qty{308.6}{K}, es decir \qty{35.4}{\celsius}, diez grados más cálido que el aire que ya estaba allí.
\end{solution}

\begin{solution}{exo:b1:fluid-envelopes-structure:6}
(a)~$T = 288.15 - 6.5 \times 5 = \qty{255.65}{K}$ y $p = \num{101325} \times
(255.65/288.15)^{5.256} = \qty{54020}{Pa}$, \qty{540.2}{hPa}. (b)~$H =
\qty{8.43}{km}$ y $p = \num{101325}\,e^{-5/8.43} = \qty{56010}{Pa}$. La presión real (estándar) es más baja: el aire allí es más frío que \qty{288.15}{K}, por tanto más denso, y la presión baja más deprisa con la altura. (c)~$p = \num{22632}\,
e^{-9/6.34} = \qty{5475}{Pa}$, el valor de
\cref{tab:b1:fluid-envelopes-structure:layers}.
\end{solution}

\begin{solution}{exo:b1:fluid-envelopes-structure:7}
(a)~$\beta\,\Delta S = \alpha\,\Delta T$ da $\Delta S = 1.67/7.61 = 0.22$.
(b)~$\Delta\rho = 1026.95 \times (-1.67\times10^{-4} \times 2 + 7.61\times10^{-4}
\times 0.3) = \qty{-0.11}{kg.m^{-3}}$: el agua se vuelve más ligera.
(c)~$\Delta\rho = 1028.11 \times (-0.53\times10^{-4} \times 2 + 7.85\times10^{-4}
\times 0.3) = \qty{+0.13}{kg.m^{-3}}$: los mismos cambios la vuelven ahora más densa, porque
el calentamiento apenas cambia la densidad del agua cercana a la congelación.
\end{solution}

\begin{solution}{exo:b1:fluid-envelopes-structure:8}
(a)~$N^2 = (9.81/260) \times (9.76 - 4)\times10^{-3} = \qty{2.17e-4}{s^{-2}}$,
$N = \qty{0.0147}{s^{-1}}$, periodo $2\pi/N = \qty{426}{s}$, unos siete
minutos. (b)~$N^2 = (9.81/270) \times 14.76\times10^{-3} = \qty{5.36e-4}{s^{-2}}$,
$N = \qty{0.0232}{s^{-1}}$, periodo \qty{271}{s}. (c)~$N$ tiende a cero y el periodo crece sin límite: la fuerza restauradora desaparece en una capa neutra.
\end{solution}

\begin{solution}{exo:b1:fluid-envelopes-structure:9}
(a)~Es más densa, $\sigma_\theta = 27.20$ frente a 26.30: el AAIW es más dulce,
pero mucho más frío. (b)~$\Delta\sigma = 1026.95 \times [-1.67\times10^{-4} \times
(4.31 - 15.43) + 7.61\times10^{-4} \times (34.30 - 35.54)] = 1.907 - 0.969 =
0.94$, frente a $0.90$ exacto. (c)~La temperatura: el enfriamiento suma
\qty{1.91}{kg.m^{-3}}, la desalinización resta \qty{0.97}{kg.m^{-3}}.
\end{solution}

\begin{solution}{exo:b1:fluid-envelopes-structure:10}
(a)~Por \cref{thm:b1:fluid-envelopes-structure:stability}, $(g/\theta)\,\mathrm
d\theta/\mathrm dz = (g/T)(\Gamma_d - \Gamma)$; $\mathrm d\theta/\mathrm dz = 0$ da $\Gamma = \Gamma_d$. (b)~$T = T_s - \Gamma_d z$, una capa de gradiente constante
$L = \Gamma_d$: por \cref{prop:b1:fluid-envelopes-structure:layer-pressure}, $p =
p_s (T/T_s)^{gM/(R\Gamma_d)}$, y $gM/(R\Gamma_d) = gMc_p/(Rg) = c_pM/R$
($= 3.5$ para aire seco). (c)~$p$ se anula donde $T = 0$, en $z = T_s/\Gamma_d =
300/9.76 = \qty{30.7}{km}$. Las capas reales de \omterm{def:b1:fluid-envelopes-structure:potential-temperature}{temperatura potencial} uniforme tienen a lo sumo unos pocos kilómetros de espesor: las tapa aire estable, como la troposfera la tapa la estratosfera.
\end{solution}

\begin{solution}{exo:b1:fluid-envelopes-structure:11}
(a)~$S = (33.787 + 36.050)/2 = 34.919$, $\theta = \qty{8.0}{\celsius}$.
(b)~Con una ecuación de estado lineal, $\sigma$ es una función afín de $S$ y de
$\theta$, así que el $\sigma$ de la mezcla es la media de las dos, 27.00.
(c)~$\alpha$ vale \qty{0.75e-4}{K^{-1}} a \qty{2}{\celsius} y
\qty{2.07e-4}{K^{-1}} a \qty{14}{\celsius}, mientras que $\beta$ cambia poco
(\num{7.80e-4} y \num{7.53e-4}). La densidad baja cada vez más deprisa al subir la temperatura: es una función cóncava de la temperatura. A lo largo de la recta de mezcla, la mezcla está a la temperatura media, donde la densidad yace por encima de la cuerda que une los dos padres, es decir, por encima de 27.00. (d)~El exceso procede de la curvatura, un término de segundo orden: crece como el cuadrado de la diferencia de temperaturas de los padres.
\end{solution}

\begin{solution}{exo:b1:fluid-envelopes-structure:12}
(a)~La diferencia de presión es $5604 - 5500 = \qty{104}{dbar}$: la partícula se enfría
$0.129 \times 0.104 = \qty{0.013}{\celsius}$, hasta \qty{0.718}{\celsius}.
(b)~El agua que la rodea a \qty{5400}{m} está a \qty{0.730}{\celsius}, más cálida: la partícula, de salinidad casi igual, es más densa que su entorno y se hunde de nuevo. La capa es estable, aunque la temperatura medida aumenta hacia abajo; su \omterm{def:b1:fluid-envelopes-structure:potential-temperature}{temperatura potencial} disminuye hacia abajo. (c)~$\theta \approx
0.7313 - 0.086 \times 5.604 = \qty{0.249}{\celsius}$, a menos de una milésima de grado del valor exacto.
\end{solution}

\begin{solution}{pb:b1:fluid-envelopes-structure:1}
\begin{enumerate}
  \item La temperatura es casi uniforme hasta \qty{50}{m}, baja deprisa hasta
        \qty{300}{m}, más despacio hasta \qty{1000}{m}, y es casi uniforme por debajo.
        La salinidad disminuye de 35.80 en la superficie a un mínimo de 34.32 a
        \qty{800}{m}, y luego aumenta hasta 34.90 a \qty{3000}{m}.
  \item $22.40 - 0.5 = \qty{21.90}{\celsius}$ se alcanza entre \qty{50}{m}
        (22.10) y \qty{75}{m} (19.20): $50 + 25 \times 0.20/2.90 = \qty{52}{m}$.
  \item $(22.10 - 13.00)/250 = \qty{0.0364}{\celsius/m}$, \qty{3.6}{\celsius}
        por \qty{100}{m}. La \omterm{def:b1:fluid-envelopes-structure:ocean-layers}{termoclina} se extiende desde la base de la \omterm{def:b1:fluid-envelopes-structure:ocean-layers}{capa de mezcla}, cerca de \qty{50}{m}, hasta unos \qty{800}{m}; es más aguda entre
        \qty{50}{m} y \qty{75}{m} (\qty{11.6}{\celsius} por \qty{100}{m}).
  \item $\sigma = 1026.95 \times [1 - 1.67\times10^{-4} \times 12.4 +
        7.61\times10^{-4} \times 0.80] - 1000 = 25.45$, \qty{0.74}{kg.m^{-3}} demasiado
        alto. La superficie está lejos del estado de referencia, y $\alpha$
        es mucho mayor a \qty{22}{\celsius} que a \qty{10}{\celsius}: la forma lineal subestima el aligeramiento por el calor.
  \item El mínimo de salinidad a \qty{800}{m}, 34.32 a \qty{4.8}{\celsius}, es
        agua intermedia antártica, formada bajo los vientos del oeste lluviosos del océano Austral. El máximo a \qty{3000}{m}, 34.90 a
        \qty{2.5}{\celsius}, es agua profunda del Atlántico Norte, formada en los mares de Groenlandia y Labrador.
  \item $(2.50 - 2.26)/3.042 = \qty{0.079}{\celsius}$ por \qty{1000}{dbar}.
  \item $\sigma_\theta$ aumenta con la profundidad: la columna es estable. La
        \omterm{def:b1:fluid-envelopes-structure:ocean-layers}{picnoclina} es donde $\sigma_\theta$ aumenta más deprisa, en los primeros cientos de metros, junto con la \omterm{def:b1:fluid-envelopes-structure:ocean-layers}{termoclina}: \num{0.82} entre
        \qty{200}{m} y \qty{800}{m}, solo \num{0.70} en los \qty{2200}{m} de debajo.
  \item $S = (35.45 + 34.32)/2 = 34.885$; $\theta = (14.90 + 4.80)/2 =
        \qty{9.85}{\celsius}$.
  \item C: 26.461; I: 27.310; mezcla: 26.886, exactamente la media de las dos (26.886), ya que la forma lineal es afín.
  \item La media de los padres es $(26.342 + 27.159)/2 = 26.751$; la mezcla, con 26.888, es \qty{0.14}{kg.m^{-3}} más densa. Esto es cabbeling.
  \item Con la forma lineal, el perfil da 26.61 a \qty{300}{m}, 26.88 a
        \qty{500}{m} y 27.31 a \qty{800}{m}: la densidad de la mezcla yace cerca de
        \qty{500}{m}. El agua allí, $(34.82, \qty{9.6}{\celsius})$, está muy
        cerca de la mezcla $(34.885, \qty{9.85}{\celsius})$: el perfil entre el agua central y el AAIW es en sí mismo casi una recta de mezcla entre ellas.
  \item $S = 0.25 \times 35.45 + 0.75 \times 34.32 = 34.60$, $\theta = 0.25
        \times 14.90 + 0.75 \times 4.80 = \qty{7.33}{\celsius}$. Interpolando el perfil entre \qty{500}{m} y \qty{800}{m} se obtiene
        $(34.57, \qty{7.2}{\celsius})$ cerca de \qty{650}{m}: casi sobre ella.
  \item En \si{K/km}: A--B 10.0; B--C 9.7; C--D $-5.5$; D--E 6.29; E--F 6.50;
        F--G 6.57; G--H 0.33; H--I $-0.50$; I--J $-0.83$.
  \item A--B inestable; B--C neutra; C--D una inversión; D--E, E--F, F--G estables;
        G--H casi isoterma, estable; H--I e I--J inversiones. La tropopausa está en G, \qty{12.5}{km} y \qty{-59.0}{\celsius}, donde la temperatura deja de disminuir.
  \item La temperatura media es \qty{214.65}{K}: $H = 287 \times 214.65/9.81 =
        \qty{6.28}{km}$, y $147\,e^{-4/6.28} = \qty{77.7}{hPa}$, coherente con los 78 medidos.
  \item El aire en A, \qty{22.0}{\celsius}, está más frío que el mar: el calor fluye del mar al aire, que se calienta desde abajo. La capa A--B, con un gradiente por encima de $\Gamma_d$, es el resultado.
  \item En B: $22.0 - 9.76 \times 0.5 = \qty{17.1}{\celsius}$, frente a 17.0; en
        C: \qty{7.4}{\celsius}, frente a 7.3; en D: \qty{3.5}{\celsius}, frente a
        9.5. La partícula es algo más cálida que el aire hasta C, y mucho más fría en D.
  \item En la base de la inversión, cerca de \qty{1.5}{km}: al entrar en la inversión, la partícula se vuelve más fría, y por tanto más densa, que el aire que la rodea, lo que detiene su ascenso.
  \item Las partículas liberadas por la capa superficial inestable suben a través de ella y la mezclan hasta que su \omterm{def:b1:fluid-envelopes-structure:potential-temperature}{temperatura potencial} es uniforme, lo que hace que su gradiente sea $\Gamma_d$; la inversión de encima detiene la mezcla.
  \item En K: A 294.15; B 293.99; C 294.08; D 300.51; E 312.65; F 329.29; G
        346.28; H 369.50; I 446.99; J 600.91.
  \item $\theta$ disminuye de A a B (inestable), es casi constante de B a C (neutra) y aumenta por encima (estable), más deprisa de C a D y por encima de G: la misma \omterm{def:b1:clastic-sediments:sorting}{clasificación} que en la pregunta~14.
  \item E--F: $(329.29 - 312.65)/4 = \qty{4.2}{K/km}$; H--I: $(446.99 -
        369.50)/4 = \qty{19.4}{K/km}$, casi cinco veces más.
  \item Con la media $\theta = \qty{408.2}{K}$: $N^2 = (9.81/408.2) \times
        0.01937 = \qty{4.66e-4}{s^{-2}}$, $N = \qty{0.0216}{s^{-1}}$, un periodo
        de \qty{291}{s}, unos cinco minutos.
  \item $N^2 = (9.81/1027) \times (26.35 - 24.79)/150 = \qty{9.9e-5}{s^{-2}}$,
        $N = \qty{0.0100}{s^{-1}}$, un periodo de \qty{630}{s}. La estratosfera inferior es la más estable, con $N^2$ casi cinco veces mayor que en la \omterm{def:b1:fluid-envelopes-structure:ocean-layers}{termoclina}.
  \item En G, $\theta = 214.15 \times (1000/186)^{2/7} = \qty{346}{K}$: la \omterm{def:b1:fluid-envelopes-structure:potential-temperature}{temperatura potencial} en la tropopausa es de unos \qty{346}{K}.
\end{enumerate}
\end{solution}
